\section{The unchanged seed's away-inertia condition}

\begin{lemma}[AC-G60-seed-scope: nonzero monodromy is outside the stated seed domain]
\label{AC-G60-seed-scope}
Let \(v\nmid\ell\) be a nonarchimedean place and let
\(r\) be a characteristic-zero \(\ell\)-adic representation.
Suppose that on an open inertia subgroup \(I_0\) it has the form
\[
 r(s)=\exp(t_\ell(s)N),\qquad N\ne0,
 \tag{G60-scope.1}
\]
where \(N\) is nilpotent and the image of \(t_\ell\) is an open
subgroup of \(\mathbb Z_\ell\).
After any finite extension of the local field, its inertia action
cannot be conjugate to a direct sum of characters.  The same conclusion
holds after any scalar character twist or extension of coefficients.
Consequently a representation with a Steinberg Weil--Deligne block
at an auxiliary place away from \(\ell\) does not satisfy the
unchanged source's \texttt{TypeAway}, and hence does not satisfy
\texttt{IsSeed.away} or \texttt{IsSeedAlt.away}.
Enlarging an auxiliary level integer does not change this assertion.
\end{lemma}
\begin{proof}
The inertia group of a finite local extension intersects \(I_0\)
in a subgroup of finite index.  Its \(t_\ell\)-image therefore
contains a nonzero element \(a\).  The matrix \(\exp(aN)\)
is unipotent and is not the identity: the finite polynomial logarithm
on unipotent matrices gives
\(\log(\exp(aN))=aN\ne0\).
A diagonalizable unipotent matrix over a characteristic-zero field
is the identity.  Thus this inertia element is not diagonalizable,
whereas every matrix in a direct sum of characters is diagonalizable
over an algebraic closure.  This is a contradiction.
Scalar multiplication by a nonzero character value preserves a
nontrivial Jordan block.  Extension of coefficients preserves both
the nonzero operator \(N\) and that block.  In the convention in which
one renormalizes the tame parameter after a ramified extension of
index \(e\), the monodromy becomes \(eN\); it is still nonzero in
characteristic zero.  Equivalently the preceding restriction argument
does not require choosing this normalization.

A Steinberg block has a nonzero nilpotent monodromy operator, so it
falls under the stated hypothesis.  The adjoint rank-three block of
a rank-two Steinberg representation does too: its monodromy is
\(\operatorname{ad}(N)\) on traceless matrices, which is nonzero
because a nonzero nilpotent rank-two matrix is not central.

Finally, the literal source definition is
\[
 \texttt{TypeAway}(r,\tau)
  =\bigl(\forall A\text{ away from }\ell,\;
          \texttt{IsDiagOn}(r,I_A,\tau)\bigr).
 \tag{G60-scope.2}
\]
There is no exception for primes dividing a level integer in this
quantifier.  Both seed records require precisely this field.
Adding the auxiliary prime to a level therefore does not remove
the contradictory inertia requirement.
\end{proof}

\paragraph{Exact source scope.}
The exact meaning of \texttt{IsDiagOn} is given in
\path{AutomorphicForms/BlochKatoSelmerVanishes.lean},
lines 33--35: one fixed conjugating matrix makes every inertia
matrix exactly diagonal with the specified characters.  Thus
this condition is stronger than a statement about the inertia
semisimplification.  The pinned seed sources are
\path{IdeleGenerators/SlopeTestRemainingChoices.lean},
\texttt{TypeAway} (lines 83--84) and \texttt{IsSeed.away}
(line 132), and
\path{ModularityArgument/ThresholdArguments/DenseThresholdCases/IsSeed.lean},
\texttt{IsSeedAlt.away} (line 31).
This excludes the proposed auxiliary-Steinberg construction from
these unchanged input records.  It does not exclude constructing a
cubic automorphic representation with a different local input or
proving a modified propagation theorem.  It is a source-domain
correction, not a closure of a main arithmetic gate.

\paragraph{Geometric dictionary.}
Restriction along a finite extension is restriction of the associated
lisse representation, in the already audited SGA~1, Expos\'e~V,
\S\S4 and 7 dictionary.  The obstruction here is the explicit
nonzero local monodromy operator; no existence assertion is obtained
from that categorical dictionary.
